\section{Project Description}

Ability-based epistemic logics are found many places in computer science, 
as arguably they provide us with a formal approach for strategic reasoning and AI planning. 
Several logics in this family have been studied lately. In particular, epistemic logics describing the notion of \textit{knowing how} received a lot of attention, 
due to their simplicity and good meta-logical properties.\par
The simplest version of \textit{knowing how}, $\mathcal{L}_{Kh}$, based on the paper by (Wang 2015\cite{Wang2015}), with the semantics reformulated in (Areces et al. 2021\cite{Areces2021}). 
This logic captures the agent's ability to guarantee the completion of a task under a given condition. 
This idea is ideal in a sense, since even though the agent manages to do something under a condition, it does not imply the agent really know how. 
For example, one can win a lottery does not imply one can guarantee how to win. The uncertainty is involved in this case. 
Therefore, another logic may be prefered, $\mathcal{L}^U_{Kh}$. In this logic, an agent may have \textit{indistinguishability} among several plans they might take, a critical concept in epistemic logic. 
Indistinguishable plans form equivalence classes, partitioning the set of available plans. Intuitively, plans considered identical by the agent are placed in the same class. 
Furthermore, this logic involves multi-agents, as opposed to a single agent. Finally, a new logic, $reg-\mathcal{L}^U_{Kh}$, a variation of the second version, 
introduces equivalence classes that are defined in terms of finite automatons. \par
In this project, we aim to model this family of knowing-how logics. 
We also intend to implement the model checking for the logic $reg-\mathcal{L}^U_{Kh}$ defined in \cite{Demri2023}.
Our project goals are defined as follows.

\begin{itemize}
    \item Get familiar with the family of knowing-how logics;
    \item Implement the the family of knowing-how logics (syntax, semantics) in Haskell;
    \item Implement testing for knowing-how languages;
    \item (Optional) Implement model checking problem for the logic $reg-\mathcal{L}^U_{Kh}$ defined in \cite{Demri2023};
    \item (Optional) Implement testing for model checking.
    \item (Optional) Implement a user-interface or at the very least an executable (using parsing) for model checking.
\end{itemize}